\chapter{TỔNG QUAN ĐỀ TÀI}
\label{chapter:tong_quan_de_tai}

\section{Bối cảnh và tính cấp thiết}
Trong tiến trình chuyển đổi số sâu rộng của ngành giáo dục, việc nghiên cứu và ứng dụng trí tuệ nhân tạo vào hoạt động dạy và học đang trở thành một xu thế tất yếu, mở ra nhiều cơ hội nâng cao chất lượng sư phạm~\cite{zawacki2019systematic, kasneci2023chatgpt, holmes2019ai, luckin2016intelligence}. Sự bùng nổ của các mô hình ngôn ngữ lớn đã đánh dấu bước phát triển vượt bậc về khả năng thấu hiểu ngữ cảnh và tạo lập ngôn ngữ tự nhiên~\cite{brown2020gpt3, zhao2023llmsurvey}. Dựa trên nền tảng này, các hệ thống trợ lý học tập có thể hỗ trợ người học tra cứu kiến thức, giải thích cặn kẽ các vấn đề khoa học phức tạp và tương tác linh hoạt theo lộ trình cá nhân hóa~\cite{kasneci2023chatgpt}. Đối với môn Khoa học tự nhiên trong Chương trình Giáo dục phổ thông cấp Trung học cơ sở, tiềm năng ứng dụng của các công nghệ này càng thể hiện rõ nét. Là môn học tích hợp liên môn giữa ba phân môn Vật lí, Hoá học và Sinh học, chương trình đòi hỏi học sinh phải tiếp thu đồng thời nhiều dạng tri thức phong phú: từ các cơ chế và quá trình biến đổi sinh học, các định luật vật lí gắn liền với công thức tính toán định lượng, cho đến bản chất các phản ứng hoá học và phương trình biến đổi chất.

Tuy vậy, việc đưa trực tiếp các mô hình ngôn ngữ lớn thương mại vào phục vụ giáo dục phổ thông hiện gặp phải nhiều rào cản kỹ thuật căn bản~\cite{kasneci2023chatgpt}. Trước hết, do được huấn luyện trên nguồn ngữ liệu mở khổng lồ từ Internet, các mô hình ngôn ngữ lớn rất dễ nảy sinh hiện tượng tạo lập thông tin sai lệch, ngụy tạo tri thức nhưng lại được diễn đạt với văn phong hết sức thuyết phục~\cite{ji2023hallucination, huang2023hallucination}, hoặc đưa ra nội dung không bám sát chuẩn kiến thức và kỹ năng của chương trình giáo dục phổ thông quốc gia. Mặt khác, tri thức nội tại của mô hình bị đóng băng tại thời điểm hoàn tất quá trình huấn luyện; các hệ thống trợ lý hội thoại thông thường hiện nay thiếu hẳn cơ chế truy xuất tường minh từ các nguồn tài liệu sư phạm chính thống như sách giáo khoa chuẩn mực hay học liệu đã qua thẩm định chuyên môn~\cite{lewis2020retrieval, gao2023ragsurvey}. Sự thiếu vắng khả năng đối chiếu và kiểm chứng nguồn gốc khiến độ tin cậy của câu trả lời bị suy giảm nghiêm trọng, gây trở ngại lớn cho việc ứng dụng công nghệ này trong môi trường giáo dục chính quy.

Để khắc phục triệt để các hạn chế trên, hướng tiếp cận sinh tăng cường truy xuất RAG~\cite{lewis2020retrieval} đã khẳng định vị thế là một giải pháp tiên tiến, kết hợp hài hòa năng lực lập luận ngôn ngữ linh hoạt của mô hình ngôn ngữ lớn với khả năng truy xuất tri thức chuẩn xác từ kho tài liệu chuyên biệt và đáng tin cậy~\cite{gao2023ragsurvey, yan2024crag, jeong2024adaptiverag}. Thay vì phụ thuộc hoàn toàn vào các trọng số tham số tĩnh, mô hình RAG chủ động tìm kiếm các đoạn thông tin liên quan nhất từ cơ sở tri thức xác thực để làm ngữ cảnh đầu vào cho mô hình ngôn ngữ tổng hợp câu trả lời hoàn chỉnh~\cite{lewis2020retrieval, guu2020realm}. Phương pháp này nâng cao vượt bậc tính chuẩn xác của câu trả lời, đảm bảo tính giải trình khoa học thông qua việc dẫn xuất tường minh nguồn gốc của từng luận điểm, đồng thời cho phép cập nhật tri thức linh hoạt mà không phát sinh chi phí tái huấn luyện mô hình tốn kém~\cite{gao2023ragsurvey}.

Đặc biệt, đặc thù nổi bật của sách giáo khoa môn Khoa học tự nhiên cấp Trung học cơ sở là nội dung lý thuyết luôn gắn liền mật thiết với hệ thống \textbf{hình ảnh trực quan, sơ đồ giải phẫu, biểu đồ thí nghiệm và phương trình khoa học}: từ cấu trúc tế bào sinh học, sơ đồ mạch điện, đồ thị chuyển động cho đến các phản ứng hoá học. Một hệ thống RAG truyền thống thuần văn bản sẽ bỏ sót toàn bộ kho tài nguyên trực quan quý giá này. Do đó, việc nghiên cứu phát triển một kiến trúc RAG \textbf{đa phương thức}, có năng lực thấu hiểu và truy xuất đồng thời cả dữ liệu văn bản lẫn hình ảnh dựa trên các mô hình liên kết thị giác -- ngôn ngữ hiện đại như CLIP~\cite{radford2021clip}, là một yêu cầu khoa học cấp thiết.

Bên cạnh bài toán mở rộng không gian tri thức đa phương thức, \textbf{thách thức trọng tâm trong nghiên cứu RAG nằm ở chất lượng của pha truy xuất}. Trong thực tế triển khai, một tỷ lệ đáng kể các đoạn tài liệu chính xác thường bị loại bỏ ở khâu chấm điểm và chọn lọc ngưỡng $top\text{-}k$, thay vì do không gian nhúng không thể tìm thấy chúng. Hệ quả là độ phủ ở quy mô vận hành thực tế thường thấp hơn nhiều so với năng lực trần lý thuyết của riêng bước tìm kiếm ngữ nghĩa nếu mở rộng số lượng đoạn ứng viên. Chính khoảng cách này phản ánh giới hạn có thể tối ưu hoá bằng phương pháp luận kiến trúc chứ không chỉ đơn thuần là tăng quy mô tham số mô hình. Cần nhấn mạnh rằng, các kết quả thực nghiệm định lượng tại Chương~\ref{chapter:thuc_nghiem_danh_gia} sẽ phân tích chi tiết mức độ cải thiện cũng như những giới hạn thực tế của khoảng cách này (xem Mục~\ref{sec:ablation}). Vì vậy, đề tài tập trung vào hai hướng nghiên cứu then chốt: (i) nâng cao độ tin cậy của luồng truy xuất đa phương thức thông qua cơ chế định tuyến ý định, truy xuất lai và bước tái xếp hạng chuyên sâu, và (ii) \textbf{kiểm chứng thực nghiệm đối chiếu khách quan}: so sánh truy xuất lai với mô hình BM25 truyền thống~\cite{robertson2009bm25}, và đối chiếu RAG đa phương thức với RAG thuần văn bản.

Xuất phát từ những yêu cầu thực tiễn và thách thức nghiên cứu nêu trên, đề tài \textbf{"Nghiên cứu phương pháp RAG đa phương thức cho trợ lý ảo thông minh hỗ trợ dạy và học môn Khoa học tự nhiên cấp Trung học cơ sở"} được thực hiện nhằm xây dựng một giải pháp trợ lý học tập thông minh, đảm bảo cung cấp học liệu số chuẩn xác, minh bạch và có khả năng truy nguyên chính xác về từng trang sách giáo khoa phục vụ học sinh và giáo viên.

\section{Mục tiêu của đề tài}
Mục tiêu bao trùm của đề tài là nghiên cứu, thiết kế, hiện thực hóa và đánh giá định lượng phương pháp RAG đa phương thức ứng dụng trong trợ lý ảo hỗ trợ dạy và học môn Khoa học tự nhiên cấp Trung học cơ sở. Các mục tiêu nghiên cứu cụ thể bao gồm:

\begin{itemize}
    \item \textbf{Mục tiêu 1 --- Xây dựng cơ chế truy vấn đa định dạng:} Xây dựng mô hình cho phép hệ thống tiếp nhận, thấu hiểu và truy xuất đồng thời trên nhiều định dạng nội dung đặc thù của môn Khoa học tự nhiên --- văn bản lý thuyết, \textbf{công thức khoa học}, hình ảnh minh hoạ, sơ đồ và biểu đồ thí nghiệm --- khắc phục triệt để hạn chế của các hệ thống thuần văn bản.
    \item \textbf{Mục tiêu 2 --- Thiết lập cơ sở tri thức vector toàn diện:} Thực hiện số hoá, tiền xử lý và lập chỉ mục \textbf{toàn bộ} chương trình môn Khoa học tự nhiên từ lớp 6 đến lớp 9 của cả ba bộ sách giáo khoa hiện hành gồm Kết nối tri thức với cuộc sống, Chân trời sáng tạo và Cánh Diều~\cite{sgk_kntt6, sgk_kntt7, sgk_kntt8, sgk_kntt9, sgk_ctst6, sgk_ctst7, sgk_ctst8, sgk_ctst9, sgk_cd6, sgk_cd7, sgk_cd8, sgk_cd9}; xử lý chuyên sâu cả văn bản lý thuyết, ký hiệu công thức và hình ảnh minh họa.
    \item \textbf{Mục tiêu 3 --- Tối ưu hóa pha truy xuất tri thức:} Hoàn thiện cơ chế định tuyến theo ý định truy vấn, phát triển mô hình truy xuất lai kết hợp giữa tìm kiếm từ khoá BM25 và tìm kiếm ngữ nghĩa vector, cùng bước tái xếp hạng bằng mô hình tương tác chéo, nhằm \textbf{thu hẹp khoảng cách} giữa độ phủ thực tế và trần độ phủ lý thuyết của hệ thống.
    \item \textbf{Mục tiêu 4 --- Thiết lập khung đánh giá thực nghiệm đối chiếu:} Xây dựng quy trình đánh giá định lượng chặt chẽ và thực hiện \textbf{thực nghiệm so sánh đối đầu} trên cùng một bộ dữ liệu kiểm thử: (i) truy xuất lai so với mô hình BM25 truyền thống, (ii) RAG đa phương thức so với RAG thuần văn bản; đồng thời tiến hành thực nghiệm bóc tách thành phần nhằm phân tích vai trò của bước tái xếp hạng và cổng lọc theo độ liên quan.
    \item \textbf{Mục tiêu 5 --- Hiện thực hóa ứng dụng trợ lý học tập trực quan:} Phát triển hoàn chỉnh ứng dụng web tương tác cho phép giáo viên và học sinh đặt câu hỏi và nhận câu trả lời kèm hình ảnh trích xuất, \textbf{công thức hiển thị chuẩn ký hiệu khoa học} và thông tin trích dẫn minh bạch đến từng trang sách giáo khoa nguồn.
    \item \textbf{Cơ sở phương pháp luận cho các mục tiêu --- Bộ dữ liệu kiểm thử chuẩn đối sánh:} Xây dựng bộ dữ liệu kiểm thử được gắn nhãn nguồn xác định theo tên sách và số trang làm chuẩn đối sánh khách quan, bao quát cả ba phân môn, bốn khối lớp và ba bộ sách giáo khoa. Pha truy xuất được lượng hoá bằng các chỉ số Precision@$k$, Recall@$k$ và MRR đối chiếu trực tiếp nhãn nguồn; pha sinh ngôn ngữ được đánh giá qua phương pháp LLM-as-a-judge~\cite{zheng2023llmjudge} với mô hình ngôn ngữ độc lập chấm điểm theo ba tiêu chí cốt lõi: Tính đúng, Độ trung thực và Độ liên quan, dựa trên các khung tiêu chuẩn đo lường RAG chuẩn hóa~\cite{es2024ragas, saadfalcon2023ares}.
\end{itemize}

\section{Đối tượng và phạm vi nghiên cứu}
\subsection{Đối tượng nghiên cứu}
\begin{itemize}
    \item Các phương pháp xử lý ngôn ngữ tự nhiên và trí tuệ nhân tạo tạo sinh.
    \item Các mô hình ngôn ngữ lớn mã nguồn mở có năng lực tuân thủ chỉ thị cao.
    \item Kiến trúc hệ thống sinh tăng cường truy xuất đa phương thức RAG xử lý đồng thời dữ liệu văn bản và hình ảnh.
    \item Cơ sở dữ liệu vector, các mô hình nhúng ngữ nghĩa và mô hình tái xếp hạng tương tác chéo.
\end{itemize}

\subsection{Phạm vi nghiên cứu}
\begin{itemize}
    \item \textbf{Về phạm vi dữ liệu:} Nghiên cứu bao quát \textbf{toàn bộ} chương trình môn Khoa học tự nhiên cấp Trung học cơ sở từ lớp 6 đến lớp 9 theo Chương trình Giáo dục phổ thông 2018, tích hợp cả ba phân môn \textbf{Vật lí, Hoá học và Sinh học}. Cơ sở tri thức được xây dựng từ 12 cuốn sách giáo khoa thuộc ba bộ sách hiện hành: Kết nối tri thức với cuộc sống~\cite{sgk_kntt6, sgk_kntt7, sgk_kntt8, sgk_kntt9}, Chân trời sáng tạo~\cite{sgk_ctst6, sgk_ctst7, sgk_ctst8, sgk_ctst9} và Cánh Diều~\cite{sgk_cd6, sgk_cd7, sgk_cd8, sgk_cd9}. Do đó, các nội dung thuộc phân môn Vật lí và Hoá học là thành phần cốt lõi trong phạm vi nghiên cứu của đề tài.
    \item \textbf{Về mặt kỹ thuật:} Hệ thống được thiết kế theo kiến trúc máy khách -- máy chủ hiện đại. Phía máy chủ chịu trách nhiệm quản lý quy trình RAG, điều phối mô hình nhúng và mô hình ngôn ngữ lớn sử dụng ngôn ngữ Python, thư viện LangChain và cơ sở dữ liệu vector ChromaDB. Phía máy khách cung cấp giao diện tương tác trực quan cho người dùng phát triển trên nền tảng Next.js.
    \item \textbf{Về đối tượng phục vụ:} Học sinh và giáo viên cấp Trung học cơ sở có nhu cầu tra cứu học liệu, ôn tập kiến thức và giải đáp thắc mắc liên quan đến chương trình môn Khoa học tự nhiên.
    \item \textbf{Giới hạn phạm vi nghiên cứu:} Đề tài không hướng đến việc tinh chỉnh toàn bộ trọng số của các mô hình ngôn ngữ lớn, mà tập trung chủ đạo vào việc hoàn thiện cấu trúc và nâng cao độ chính xác của quy trình truy xuất tri thức cùng khả năng tích hợp đa phương thức.
\end{itemize}

\section{Cấu trúc của báo cáo đồ án}
Báo cáo đồ án tốt nghiệp được tổ chức thành 5 chương với nội dung tóm tắt như sau:

\begin{itemize}
    \item \textbf{Chương 1: Tổng quan đề tài.} Trình bày bối cảnh thực tiễn, tính cấp thiết, mục tiêu nghiên cứu, đối tượng, phạm vi và bố cục tổng thể của báo cáo.
    \item \textbf{Chương 2: Cơ sở lý thuyết.} Hệ thống hóa các kiến thức nền tảng phục vụ nghiên cứu, bao gồm các nguyên lý xử lý ngôn ngữ tự nhiên, mô hình ngôn ngữ lớn, kiến trúc RAG, hệ quản trị cơ sở dữ liệu vector, kỹ thuật biểu diễn đa phương thức và các phương pháp truy xuất kết hợp tái xếp hạng.
    \item \textbf{Chương 3: Phương pháp thực hiện và kiến trúc hệ thống.} Mô tả chi tiết kiến trúc tổng thể của hệ thống trợ lý ảo và các công nghệ nền tảng. Trình bày quy trình trích xuất và tiền xử lý dữ liệu cho văn bản, công thức và hình ảnh; cách thức thiết kế luồng RAG đa phương thức và giải pháp giao diện người dùng.
    \item \textbf{Chương 4: Thực nghiệm, đánh giá và thảo luận.} Trình bày môi trường thực nghiệm, quy trình xây dựng bộ dữ liệu kiểm thử gắn nhãn nguồn khách quan, hệ thống chỉ số đánh giá hai pha truy xuất và sinh ngôn ngữ; phân tích chi tiết kết quả định lượng qua các bảng đối chiếu thực nghiệm, thảo luận chuyên sâu về các phát hiện khoa học và nhận diện các giới hạn kỹ thuật của hệ thống.
    \item \textbf{Chương 5: Kết luận và hướng phát triển.} Tổng kết các kết quả nghiên cứu và mức độ hoàn thành so với các mục tiêu đề ra, đồng thời đề xuất lộ trình và giải pháp phát triển hệ thống trong tương lai.
\end{itemize}
